\section{更广泛的影响}
\label{app:broader-impact}

本工作提出了\emph{agentic harness engineering}（AHE），这是一个自演化闭环：coding agent 依据执行反馈来编辑自己的脚手架（系统提示词、工具、中间件与长期记忆）。AHE 降低了打造具有竞争力的 coding agent 所需的人力工程成本：没有专职 harness 团队的实践者，也能在同一基座模型上获得比仅提示词的自演化基线更高的 pass@1，且单次试验的 token 成本更低，这让学术团队、中小型组织与教育场景都更容易用上有能力的编程助手。由于该闭环把反复出现的协同模式沉淀到工具、中间件与记忆中，而不是沉淀到越来越长的提示词里，它降低了单次调用的算力开销以及推理阶段相应的能耗足迹。它还使得 harness 能够快速迭代，从而跟上基座模型发布的节奏，让周边脚手架不必在每个新模型面前都落后一步。负面的社会影响以及相应的泛化风险在 \S\ref{sec:limitations} 中讨论。
